电磁相互作用是自然界的基本相互作用之一。与质点力学相比，静电学引入了一个新的描述层次：
带电体先改变周围空间的物理性质，形成电场；另一个电荷进入该区域后，再受到电场的作用。
本讲从电荷与库仑定律出发，建立电场强度、电场线和电通量的概念，最后给出高斯定律的积分
形式与微分形式。

\subsection*{物质结构与电荷}

宏观物质由原子、分子组成。原子核带正电，核外电子带负电。物体所带电荷量记为 $q$，
国际单位为库仑（\si{C}）。电荷具有以下基本性质：
\begin{enumerate}[label=\arabic*.,leftmargin=2em]
  \item \textbf{可加性：}复合系统的总电荷等于各部分电荷的代数和；
  \item \textbf{守恒性：}孤立系统内的总电荷在任何过程中保持不变；
  \item \textbf{量子化：}宏观物体的电荷量是元电荷
  \begin{equation}
    e=1.602176634\times10^{-19}\ \si{C}
  \end{equation}
  的整数倍。电子带电 $-e$，质子带电 $+e$。
\end{enumerate}

\matterStructureDiagram

点电荷是对带电体的理想化：当带电体的线度远小于问题中的距离尺度时，可以忽略其形状和
大小，只保留位置与总电荷量。

\subsection*{库仑定律}

真空中两个静止点电荷 $q_1,q_2$ 相距 $r$，它们之间的静电力为
\begin{equation}
  \boxed{\vect F_{12}=
  \frac{1}{4\pi\varepsilon_0}\frac{q_1q_2}{r^3}\vect r_{12}},
\end{equation}
其中 $\vect r_{12}$ 从 $q_1$ 指向 $q_2$，$\varepsilon_0$ 为真空介电常数，并有
\begin{equation}
  k\equiv\frac{1}{4\pi\varepsilon_0}
  \approx8.99\times10^9\ \si{N.m^2.C^{-2}}.
\end{equation}
同号电荷相斥，异号电荷相吸。多个点电荷同时存在时，总静电力满足矢量叠加原理。

\coulombDiagram

Coulomb 用扭秤把很小的静电力转化为可测的扭转角。若悬丝的扭转系数为 $\kappa$，电力对悬丝
轴的力臂为 $d$，平衡时满足
\begin{equation}Fd=\kappa\theta.\end{equation}
保持电荷量不变并改变两带电小球的距离，测得的 $\theta$ 就与 $F$ 成正比；实验所得
$\theta\propto r^{-2}$，从而给出平方反比关系。下图右侧装置中的细丝、横杆和带电小球正是
这一测量链条的实物结构。

\coulombHistoryContextDiagram

库仑力与万有引力同为平方反比中心力，许多数学结论可以类比使用；但电荷有正、负两种，
因而静电力既可能吸引，也可能排斥。两个点电荷的相互作用势能为
\begin{equation}
  U=\frac{1}{4\pi\varepsilon_0}\frac{q_1q_2}{r},
\end{equation}
零势能取在两电荷相距无穷远处。

\subsection*{电场与电场线}

电场强度描述电场在空间某点对电荷的作用。把电荷量为 $q$ 的试探电荷置于该点，若它受到
电场力 $\vect F$，则
\begin{equation}
  \boxed{\vect E=\frac{\vect F}{q}},
  \qquad \vect F=q\vect E.
\end{equation}
电场由场源电荷决定，与试探电荷的大小和正负无关。点电荷 $Q$ 在位置矢量 $\vect r$ 处产生
的电场为
\begin{equation}
  \boxed{\vect E(\vect r)=
  \frac{1}{4\pi\varepsilon_0}\frac{Q}{r^3}\vect r}.
\end{equation}
对于离散电荷系与连续电荷分布，叠加原理分别给出
\begin{equation}
  \vect E=\frac{1}{4\pi\varepsilon_0}
  \sum_i\frac{q_i\vect r_i}{r_i^3},
  \qquad
  \vect E=\frac{1}{4\pi\varepsilon_0}
  \int\frac{\vect r}{r^3}\,\dd q.
\end{equation}
式中的 $\vect r_i$ 或 $\vect r$ 均从电荷元指向场点。

电场线是在每一点都与当地电场方向相切的曲线。它不是空间中真实存在的细线，也不代表
带电粒子的运动轨迹。静电场的电场线从正电荷或无穷远出发，终止于负电荷或无穷远；不同
电场线不能相交，否则交点处的电场方向将不唯一。

\fieldLineDiagram

若在平面内用参数方程 $\vect r(s)$ 描述电场线，它满足
\begin{equation}
  \frac{\dd\vect r}{\dd s}\parallel\vect E,
\end{equation}
在直角坐标中可写为
\begin{equation}
  \frac{\dd x}{E_x}=\frac{\dd y}{E_y}=\frac{\dd z}{E_z}.
\end{equation}

\subsection*{电通量与高斯定律}

匀强电场 $\vect E$ 穿过面积为 $S$ 的平面，取该面的单位法向量为 $\vect n$，定义面积矢量
$\vect S=S\vect n$，则电通量为
\begin{equation}
  \Phi_E=\vect E\cdot\vect S=ES\cos\theta.
\end{equation}

\fluxDiagram

对一般曲面，把曲面分割为无限小面积元 $\dd\vect S$，电通量为
\begin{equation}
  \Phi_E=\int_S\vect E\cdot\dd\vect S.
\end{equation}
闭合曲面的面积元方向规定为由内向外。真空中的高斯定律为
\begin{equation}
  \boxed{\oint_S\vect E\cdot\dd\vect S
  =\frac{Q_{\rm in}}{\varepsilon_0}},
\end{equation}
其中 $Q_{\rm in}$ 是闭合曲面包围的净电荷。面外电荷会改变曲面各处的电场，却不改变穿出
整个闭合面的净通量。

\gaussDiagram

\gaussContextDiagram

从左至右看，肖像后的第一幅小图对应面积元通量
$\dd\Phi_E=\vect E\cdot\dd\vect S=E\dd S\cos\theta$；后两幅图比较点电荷位于闭合面内、外
两种情况。电荷在面内时，无论闭合面怎样变形，包围同一电荷的总通量仍为
$q/\varepsilon_0$；电荷在面外时，电场线穿入与穿出闭合面的贡献相互抵消，净通量为零。
这正对应高斯定律中真正需要判断的量 $Q_{\rm in}$，而不是曲面的形状或面外电荷的多少。

高斯定律在任何电荷分布下都成立，但只有当电场具有足够高的对称性时，它才适合直接计算
电场。例如球对称、柱对称和平面对称系统中，可选择与对称性匹配的高斯面，使积分中的
$E$ 成为常量或使部分面积上的通量为零。

利用散度定理
\begin{equation}
  \oint_S\vect E\cdot\dd\vect S
  =\iiint_V(\nabla\cdot\vect E)\,\dd V
\end{equation}
以及 $Q_{\rm in}=\iiint_V\rho\,\dd V$，可得高斯定律的微分形式
\begin{equation}
  \boxed{\nabla\cdot\vect E=\frac{\rho}{\varepsilon_0}}.
\end{equation}
它表明电荷密度是静电场的“源”：正电荷附近有净电通量流出，负电荷附近有净电通量流入。

\gaussDifferentialCubeDiagram

\subsection*{电荷分布的三种连续模型}

当大量微小带电粒子密集分布时，逐个求和既不方便也没有必要。按电荷所占据的几何维数，常用
线电荷密度、面电荷密度和体电荷密度：
\begin{equation}
 \lambda=\frac{\dd q}{\dd l},
 \qquad
 \sigma=\frac{\dd q}{\dd S},
 \qquad
 \rho=\frac{\dd q}{\dd V}.
\end{equation}
总电荷分别由 $Q=\int\lambda\,\dd l$、$Q=\int\sigma\,\dd S$、
$Q=\int\rho\,\dd V$ 得到。选用哪一种模型，取决于带电体厚度与观察尺度的比较；同一根有有限
半径的带电细棒，在远处可近似为线电荷，在极远处又可进一步近似为点电荷。

对连续分布，场点记为 $\vect r$，源点记为 $\vect r'$，电场的统一表达式为
\begin{equation}
 \vect E(\vect r)=\frac1{4\pi\varepsilon_0}
 \int\frac{\vect r-\vect r'}{|\vect r-\vect r'|^3}\,\dd q(\vect r').
\end{equation}
必须区分场点坐标和源点坐标；分母是二者距离的三次方，分子提供从源点指向场点的方向。对称性
可以在积分前消去部分分量。例如场点位于均匀带电圆环轴线上时，所有横向分量成对抵消，只需
积分轴向分量。

叠加原理来自静电场方程的线性。若电荷分布写成几部分之和，电场就是各部分电场的矢量和。
“挖去空腔”常用负电荷叠加：先填满整个带电体，再在空腔位置叠加一个形状相同、密度相反的
带电体。这样既保留原问题的几何条件，又能直接使用已知的高对称场。

\subsection*{电场线、试探电荷与粒子轨迹}

试探电荷应足够小，使它对原有电荷分布的扰动可以忽略。定义式 $\vect E=\vect F/q$ 中的 $q$
带有符号，因此无论试探电荷取正取负，所得 $\vect E$ 都相同。教学图中通常想象用正试探电荷
确定电场方向；负电荷受到的力与电场方向相反。

电场线的疏密没有绝对数值含义，只有在同一幅图采用统一画法时，才能定性比较场强。给定一小
面积元，穿过它的电场线数可被安排为与电通量成正比；因此电场线不会在无电荷区域凭空开始或
终止。静电场线也不能形成闭合曲线，因为若沿闭合电场线前进，$\vect E\cdot\dd\vect l$ 始终
同号，会与静电场环流为零矛盾。

带电粒子的轨迹由 Newton 第二定律决定：
\begin{equation}
 m\ddot{\vect r}=q\vect E(\vect r).
\end{equation}
轨迹方向由粒子的初速度和电荷符号共同决定，一般不会与电场线重合。只有正电荷从静止释放，且
运动过程中电场方向不发生横向改变等特殊情况下，轨迹才可能沿电场线。

\subsection*{高斯定律怎样真正用于求场}

高斯定律给出的是闭合面上法向电场的积分，单靠它不能确定曲面每一点的电场。要把积分化成
$ES$，必须同时证明：电场在所选曲面的有效部分大小相等，并且方向与面积元平行或垂直。常见的
三类选择是：
\begin{itemize}
  \item 球对称分布选同心球面，$\vect E$ 沿径向且同一球面上大小相等；
  \item 无限长柱对称分布选同轴圆柱面，侧面贡献 $E(2\pi rl)$，两个端面通量为零；
  \item 无限大平面对称分布选跨过平面的薄柱，两个端面贡献相等，侧面通量为零。
\end{itemize}

以半径 $R$、体电荷密度均匀为 $\rho$ 的实心球为例。球内半径 $r<R$ 的高斯面包围电荷
$Q_r=4\pi r^3\rho/3$，故
\begin{equation}
 E(r)=\frac{\rho r}{3\varepsilon_0}.
\end{equation}
球外则只看总电荷 $Q=4\pi R^3\rho/3$，有
\begin{equation}
 E(r)=\frac{\rho R^3}{3\varepsilon_0r^2}.
\end{equation}
内外结果在 $r=R$ 连续。若表面另有面电荷，法向电场会在表面发生跃变；把跨越表面的薄柱作为
高斯面可得
\begin{equation}
 E_n^{(+)}-E_n^{(-)}=\frac{\sigma}{\varepsilon_0}.
\end{equation}

静电条件 $\nabla\times\vect E=0$ 还给出切向分量连续：
$E_t^{(+)}=E_t^{(-)}$。因此一张理想面电荷会使法向分量跳变，却不会使切向分量跳变。这两个
边界关系在导体、电介质和分区电场问题中会反复使用。

\subsection*{积分形式与微分形式的互补}

积分形式关注一个有限闭合面包围的总电荷，适合利用整体对称性；微分形式关注一点附近的电场
散度，适合研究局部变化。二者完全等价，但并不意味着“某点散度为零就有电场为零”。在没有
电荷的区域，$\nabla\cdot\vect E=0$，点电荷产生的非零电场仍可以穿过该区域，只是任何不包围
电荷的小闭合面的净通量为零。

静电场还满足
\begin{equation}
 \nabla\times\vect E=\vect0,
 \qquad
 \oint\vect E\cdot\dd\vect l=0.
\end{equation}
高斯定律决定电场的“源”，零旋度条件决定它没有静电环流。两条性质和适当边界条件合在一起，
才能唯一确定静电场。

\subsection*{电荷、叠加原理与连续分布}

普通物质由原子组成，原子核带正电、电子带负电。宏观带电通常是电子转移的结果，而不是产生了
新的电荷。电荷具有代数可加性并满足守恒；孤立系统总电荷不随内部过程改变。微观上电荷以元电荷
$e$ 为基本单位量子化，质子带 $+e$、电子带 $-e$。宏观物体所含带电粒子数极多，才常把电荷分布
近似为连续函数。

真空中两个静止点电荷的相互作用满足
\[
 \vec F_{12}=\frac1{4\pi\varepsilon_0}
 \frac{q_1q_2}{r_{12}^3}\vec r_{12}.
\]
矢量形式同时包含大小和方向，同号相斥、异号相吸。多个电荷共同作用时，先分别求每个电荷产生
的力或电场，再作矢量和，这就是线性叠加原理。连续分布则把求和变为积分：
\[
 \vec E(\vec r)=\frac1{4\pi\varepsilon_0}
 \int\frac{\vec r-\vec r'}{|\vec r-\vec r'|^3}\,\dd q'.
\]
选择线、面或体电荷元并不是形式问题，而由带电体几何与观察尺度决定。

\subsection*{场线、电通量与高斯定律}

电场定义为单位正试探电荷所受力，$\vec E=\vec F/q$。试探电荷应足够小，避免显著改变原来的
电荷分布。电场线在每点的切线方向给出 $\vec E$ 方向，疏密只定性表示强弱；静电场线从正电荷
或无穷远出发，终止于负电荷或无穷远，彼此不相交，也不形成闭合曲线。场线是一种图示工具，
并非真实物质轨迹。

通过有向面元的电通量为
\[
 \dd\Phi_E=\vec E\cdot\dd\vec S,
 \qquad
 \Phi_E=\oint_S\vec E\cdot\dd\vec S.
\]
闭合面的面元规定向外。高斯定律
\[
 \oint_S\vec E\cdot\dd\vec S=\frac{Q_{\rm in}}{\varepsilon_0}
\]
对任意闭合面成立，但只有电荷分布具有球、柱或平面对称时，才能把场强从积分号中直接提出。
面外电荷虽不贡献净通量，却仍可能在面上产生非零电场；“净通量为零”也不等于“处处无场”。

球对称分布外部的场只由总电荷决定，等效为点电荷；均匀带电实心球内部包围电荷随 $r^3$ 增长，
故场强随 $r$ 线性增长。无限长均匀线电荷有
$E=\lambda/(2\pi\varepsilon_0r)$，无限大均匀薄板单侧有
$E=\sigma/(2\varepsilon_0)$。有限物体只能在离边缘足够远或观察距离满足相应条件时近似使用这些
结果，不能把理想对称结论无条件搬到边缘附近。

偶极子由相距很近的异号等量电荷构成，偶极矩 $\vec p=q\vec d$ 从负电荷指向正电荷。远处电场
按 $1/r^3$ 衰减，比点电荷的 $1/r^2$ 更快；总电荷为零并不意味着场处处为零，只是最低阶的单极项
消失。更复杂电荷分布可依次用单极、偶极、四极等多极矩描述其远场。


\subsection*{从物质结构到宏观电荷}

电与磁最初作为不同现象被研究，电磁感应的发现揭示了二者之间的联系，Maxwell 方程组又把
电场、磁场和光统一在同一理论中。静电学只处理电荷分布不随时间改变的情形，却是理解完整电磁
理论的起点。宏观物体是否带电，最终由内部正、负电荷的数量差决定；物体的尺寸和质量远大于
电子、质子等微观粒子，宏观电量却仍表现为元电荷的整数倍。

物质具有层级结构。原子由原子核和核外电子组成，原子核又由质子、中子构成，核子内部还存在
夸克与胶子等更深层结构。不同层级的典型长度相差许多数量级。讨论厘米尺度带电体时，不必追踪
每一个微观粒子，只需用总电荷、线电荷密度、面电荷密度或体电荷密度描述；讨论原子内部时，
连续分布近似便可能失效。图29.1(a)给出的长度标度用来说明模型选择，而不是把微观对象按实际比例
塞在同一张图中。

电荷守恒比“电子不会消失”更一般。成对产生、湮灭、核反应或粒子转化可以改变带电粒子的种类
和数目，但孤立系统的总电荷不变。电荷可加性意味着复合系统电量是各部分代数和；量子化则说明
通常可写 $Q=Ne$。在夸克层次会出现分数电荷，但自由宏观物体的净电荷仍以元电荷为基本尺度。

\subsection*{Coulomb 定律的适用范围与场的引入}

Coulomb 定律最初由扭秤实验定量建立。横杆绕悬丝转动，外部带电小球对杆端小球施加近似切向
的静电力，力矩 $Fd$ 与悬丝回复力矩 $\kappa\theta$ 平衡。如果把外电荷画在杆端小球的径向方向，
并把力也画成沿杆方向，力矩便会变为零，实验原理就被破坏。因此装置图中的电荷位置、力方向和
力臂必须同时满足几何关系。

点电荷模型要求带电体尺寸远小于相互距离。对球对称电荷分布，球外场可以等效为球心点电荷；
对一般形状则只能在足够远处作近似。Coulomb 定律还是静电规律：当电荷高速运动或分布随时间
变化时，磁效应和传播延迟不能忽略，不能继续把相互作用理解成瞬时的平方反比力。

引入电场以后，相互作用被分成两步：场源电荷建立空间中的场，试探电荷在当地电场中受力。场是
空间位置和时间的函数，带电粒子受到的力只由它所在位置的场决定。静止点电荷的场与 Coulomb
定律完全一致；对于变化电荷分布，场的描述仍然成立，而直接的“超距作用”图像不再可靠。

电场线只编码场的方向和相对强弱。静电场线不能相交、不能在无电荷处突然中断，也不能形成闭合
回路。振荡电荷产生的是随时间传播的电磁场，某一时刻的线形会随时间改变；因此原讲义把振荡
电荷与静态同号、异号电荷并列，是为了提醒读者电场概念并不限于静电场，而不是把三幅图都当作
时间不变的静电场线。

\subsection*{Gauss 定律的积分与微分图像}

闭合面的电通量是所有面积元上 $\vect E\cdot\dd\vect S$ 的代数和。面积元法向规定向外，场线
穿出给正贡献，穿入给负贡献。面外电荷的场线若穿过闭合面，必然一处进入、另一处离开，净贡献
为零；面内正电荷对应净流出，面内负电荷对应净流入。因此
\[
  \oint_S\vect E\cdot\dd\vect S=\frac{Q_{\rm in}}{\varepsilon_0}
\]
只取决于面内净电荷，与闭合面的形状以及面外电荷的位置无关。

高斯定律永远成立，但不总能方便地求场。只有球对称、无限柱对称或无限平面对称等情形，才能
在所选高斯面上把 $E$ 从积分中提出，或确认某些面上的通量为零。把“高斯定律成立”与“可用
高斯定律直接求 $E$”混为一谈，是静电题最常见的逻辑错误之一。

微分形式可以由一个边长分别为 $\Delta x,\Delta y,\Delta z$ 的小长方体理解。垂直于 $x$ 轴的
两个面，净通量在一阶近似下为
\[
  \frac{\partial E_x}{\partial x}\Delta x\Delta y\Delta z.
\]
另外两对表面分别给出 $\partial E_y/\partial y$ 和 $\partial E_z/\partial z$ 的贡献。六面相加，
再除以体积并令尺寸趋于零，便得到
\[
  \nabla\cdot\vect E=\frac{\rho}{\varepsilon_0}.
\]
图29.5(a)的立方体不仅是装饰，它明确了相对表面、外法向和差分位置；缺少这些位置关系，读者很难
从积分定律真正过渡到局部微分定律。

微分方程给出的是某一点附近电场的“源强”，并不能单独决定整个空间的场。相同的
$\nabla\cdot\vect E=\rho/\varepsilon_0$，配上不同边界条件会得到不同解；静电场还同时满足
$\nabla\times\vect E=0$。因此求实际电场时，需要把局部方程、电荷分布、导体表面条件和无穷远
条件放在一起。积分形式适合利用整体对称性，微分形式适合研究局部变化，二者是同一定律的不同
表达，而不是两条互不相干的公式。

跨过面电荷时，电场法向分量发生跃变
$E_{2n}-E_{1n}=\sigma/\varepsilon_0$，切向分量在静电条件下连续。这个边界关系可以由薄柱
高斯面直接得到，也是检查导体、带电薄膜和分区电场图是否画对的快速方法。

\subsection*{叠加原理与模型边界}

电场方程对场强是线性的，因此多个电荷分布产生的总场等于各部分单独产生的场之和。利用这一点，
复杂分布常被拆成“完整对称分布加上负的被挖去部分”，球形空腔、非均匀带电体和有限厚板都可
这样处理。叠加的是电场矢量；场强大小一般不能直接相加，必须先统一坐标分量。

“点电荷”“无限长直线”“无限大平面”都是带适用范围的模型。离带电体足够远时可用点电荷近似；
到直线中段的距离远小于线长时可忽略端点；观察尺度远小于平面尺寸时才可使用无限平面结果。若
题目进入端部、边缘或近场区域，对称性被破坏，高斯定律仍然成立，但不能再把场强当作高斯面上的
常量。先判断尺度关系，再选择模型，是使用高斯定律之前不可省略的一步。


\Needspace{0.45\textheight}
\begin{origexamplebox}{例题 29.1【统统算出来】}
利用积分或高斯定律计算下列均匀带电系统产生的电场：
\begin{enumerate}[label=(\arabic*),leftmargin=2.2em]
  \item 无限长直导线外距导线 $a$ 处，线电荷密度为 $\lambda$；
  \item 有限长直导线中垂线外距中点 $a$ 处，两端对中垂线所张角分别为 $\theta_1,\theta_2$；
  \item 无限大均匀带电平面外，面电荷密度为 $\sigma$；
  \item 仅一个方向无限延伸的均匀带电条带外距条带 $a$ 处，两侧边缘所张角为
  $\theta_1,\theta_2$；
  \item 半径为 $R$、面电荷密度为 $\sigma$ 的圆盘轴线上，距圆盘中心 $a$ 处。
\end{enumerate}
\finiteLineDiagram
\end{origexamplebox}

\begin{xesolution}
记 $k=1/(4\pi\varepsilon_0)$。

\textbf{(1) 无限长直导线。}取与导线同轴、半径 $a$、长度 $L$ 的圆柱面。电场沿径向，
且仅圆柱侧面有通量。由高斯定律
\[
  E(2\pi aL)=\frac{\lambda L}{\varepsilon_0},
\]
故
\[
  \boxed{E=\frac{\lambda}{2\pi\varepsilon_0a}}.
\]
当 $\lambda>0$ 时方向背离导线，$\lambda<0$ 时方向指向导线。
\lineFieldDiagram

\textbf{(2) 有限长直导线。}以中垂线指向场点为法向正方向，沿导线取坐标 $x$，电荷元
$\dd q=\lambda\dd x$ 到场点的距离为 $r=\sqrt{a^2+x^2}$。它产生的场元为
\[
  \dd E=\frac{k\lambda\dd x}{a^2+x^2}.
\]
将场元分解到法向和导线方向，并令 $x=a\tan\theta$，得到
\[
  \boxed{E_\perp=\frac{k\lambda}{a}
  (\sin\theta_1+\sin\theta_2)},
\]
\[
  \boxed{E_\parallel=\frac{k\lambda}{a}
  (\cos\theta_1-\cos\theta_2)}.
\]
$E_\parallel$ 的正方向取向张角为 $\theta_1$ 的一端；若线段关于中垂线对称，
$\theta_1=\theta_2=\theta$，则切向分量抵消，
\[
  \boxed{E=\frac{2k\lambda}{a}\sin\theta}.
\]
令 $\theta\to\pi/2$，即回到无限长直导线的结果。

\textbf{(3) 无限大平面。}取跨过平面的薄圆柱高斯面。由平面对称性，两底面的电场大小
相同且垂直平面，侧面无通量，因而
\[
  2ES=\frac{\sigma S}{\varepsilon_0},
  \qquad
  \boxed{E=\frac{\sigma}{2\varepsilon_0}}.
\]

\textbf{(4) 无限长条带。}把条带分成许多与无限方向平行的细线，每条细线的线电荷密度
$\dd\lambda=\sigma\dd x$。设场点到条带所在平面的垂足为坐标原点，条带两边分别位于
$x=-a\tan\theta_2$ 和 $x=a\tan\theta_1$。叠加无限长直线的电场，有
\[
  \boxed{E_\perp=\frac{\sigma}{2\pi\varepsilon_0}
  (\theta_1+\theta_2)},
\]
\[
  \boxed{E_\parallel=
  \frac{\sigma}{2\pi\varepsilon_0}
  \ln\frac{\cos\theta_1}{\cos\theta_2}}.
\]
这里 $E_\parallel$ 的正方向取向张角为 $\theta_1$ 的一侧。
若场点位于条带中垂面，$\theta_1=\theta_2=\theta$，则
\[
  \boxed{E=\frac{\sigma\theta}{\pi\varepsilon_0}}.
\]
当条带宽度趋于无穷，$\theta\to\pi/2$，再次得到 $E=\sigma/(2\varepsilon_0)$。

\textbf{(5) 均匀带电圆盘。}把圆盘分成半径 $s$、宽度 $\dd s$ 的同心圆环，
$\dd q=2\pi\sigma s\dd s$。圆环在轴线场点处的横向分量相消，轴向场元为
\[
  \dd E_x=k\frac{a\,\dd q}{(a^2+s^2)^{3/2}}.
\]
从 $s=0$ 积分到 $R$，得到
\[
  \boxed{E_x=\frac{\sigma}{2\varepsilon_0}
  \left(1-\frac{a}{\sqrt{a^2+R^2}}\right)}.
\]
若圆盘边缘对场点的半张角为 $\theta$，则
\[
  \boxed{E_x=\frac{\sigma}{2\varepsilon_0}(1-\cos\theta)}.
\]
\diskFieldDiagram
\end{xesolution}

\Needspace{0.35\textheight}
\begin{origexamplebox}{例题 29.2【奇妙的电场线】}
点电荷 $+Q$、$-q$ 分别位于 $(-a,0,0)$、$(a,0,0)$，且 $Q>q>0$。系统关于两电荷连线
旋转对称。回答：
\begin{enumerate}[label=(\arabic*),leftmargin=2.2em]
  \item $+Q$ 发出的电场线能否到达无穷远？$-q$ 接收的电场线是否来自无穷远？
  \item 空间中是否存在电场为零的点？该点附近的电场线有何特征？
  \item 无穷远处的电场有什么特征？
  \item 写出电场线方程的一种隐式表示。
\end{enumerate}
\unequalChargeDiagram
\end{origexamplebox}

\begin{xesolution}
\textbf{(1)} 电场线条数只能用来表示通量比例，不是绝对数量。由于净电荷为 $Q-q>0$，
$+Q$ 发出的电通量中有一部分终止于 $-q$，其余必然延伸到无穷远。$-q$ 所接收的全部电场线
都可来自 $+Q$，不需要由无穷远进入。

\textbf{(2)} 零场点只能在两电荷连线延长线上。两电荷之间的电场同向，不可能抵消；
在 $+Q$ 左侧，大电荷又更近，也不能抵消。因此零场点在 $-q$ 的右侧。设其坐标为 $x_0>a$，
则
\[
  \frac{kQ}{(x_0+a)^2}=\frac{kq}{(x_0-a)^2},
\]
解得
\[
  \boxed{x_0=a\frac{\sqrt Q+\sqrt q}{\sqrt Q-\sqrt q}}.
\]
零场点附近的场具有鞍点结构：沿两电荷连线方向，电场线从零场点两侧离开；沿垂直方向，
电场线向该轴线汇聚。电场线本身不穿过电场为零的点，因为该点没有确定的切向方向。

\textbf{(3)} 当 $r\gg a$ 时，观察者分辨不出两电荷的间距，最低阶电场等效于位于原点的
净电荷 $Q-q$：
\[
  \boxed{\vect E\simeq k\frac{Q-q}{r^2}\uvec r}.
\]
因此无穷远处的电场线近似沿径向向外，偶极修正比主项多衰减一个 $1/r$。

\textbf{(4)} 在包含两电荷的任一子午面内，设场点相对于 $+Q$、$-q$ 的极角分别为
$\theta_+$、$\theta_-$，二者的极轴都取沿 $+x$ 方向。由轴对称电场的通量函数可得，
每一条电场线满足
\[
  \boxed{Q(1-\cos\theta_+)-q(1-\cos\theta_-)=C},
\]
其中 $C$ 由所选电场线决定。等价地，也可写成
\[
  Q\cos\theta_+-q\cos\theta_-=C'.
\]
这是一种隐式方程；绕 $x$ 轴旋转即可得到三维电场线族。
\end{xesolution}

\Needspace{0.37\textheight}
\begin{origexamplebox}{例题 29.3【球壳空洞】}
一个球体内部挖去偏心球形空腔，未挖去部分具有均匀体电荷密度 $\rho$。
\begin{enumerate}[label=(\arabic*),leftmargin=2.2em]
  \item 求空腔内的电场；
  \item 分析它与万有引力中同类模型的联系，并讨论题目若把球体称为“金属球体”会有什么问题。
\end{enumerate}
\cavityDiagram
\end{origexamplebox}

\begin{xesolution}
设完整大球球心为 $O$，空腔球心为 $C$，$\vect d=\overrightarrow{OC}$。把实际电荷分布分解为：
以 $O$ 为球心、密度为 $\rho$ 的完整均匀带电球，加上以 $C$ 为球心、密度为 $-\rho$ 的小球。
均匀带电实心球内部距球心位置矢量为 $\vect r$ 的点，电场为
\[
  \vect E=\frac{\rho}{3\varepsilon_0}\vect r.
\]
对空腔内任一点 $P$，记 $\vect r=\overrightarrow{OP}$、$\vect r'=\overrightarrow{CP}$，
则 $\vect r=\vect d+\vect r'$。叠加两球的电场：
\[
  \vect E=\frac{\rho}{3\varepsilon_0}\vect r
  -\frac{\rho}{3\varepsilon_0}\vect r'
  =\boxed{\frac{\rho}{3\varepsilon_0}\vect d}.
\]
所以空腔内电场是匀强电场，方向由大球球心指向空腔球心。

万有引力也服从平方反比定律，均匀球形空腔可用完全相同的叠加方法计算，只需把电学常量与
质量密度替换为引力常量与质量密度；引力只有吸引，最终场方向相应改变。

若把球体称为“金属球体”并同时规定内部均匀体电荷密度，这在静电平衡下不成立。导体内
自由电荷会重新分布，导体材料内部电场必须为零。因此上述答案适用于均匀体电荷被冻结的
绝缘球体；若确为处于静电平衡的导体，必须按导体边界条件重新求解。
\end{xesolution}

\Needspace{0.35\textheight}
\begin{origexamplebox}{例题 29.4【电偶极子】}
一正一负两个点电荷的电荷量分别为 $+q$、$-q$，间距为 $l$。求 $r\gg l$ 处电场强度的
大小和方向，并求在该近似下的电场线方程。
\dipoleDiagram
\end{origexamplebox}

\begin{xesolution}
定义电偶极矩
\[
  \vect p=q\vect l,
\]
方向由负电荷指向正电荷。取偶极子中心为极点，以 $\vect p$ 方向为极轴。当 $r\gg l$ 时，
两点电荷势的最低阶非零项为
\[
  \varphi(r,\theta)=k\frac{p\cos\theta}{r^2}.
\]
由 $\vect E=-\nabla\varphi$ 得
\[
  \boxed{E_r=\frac{2kp\cos\theta}{r^3}},
  \qquad
  \boxed{E_\theta=\frac{kp\sin\theta}{r^3}}.
\]
因此
\[
  \boxed{E=\frac{kp}{r^3}\sqrt{1+3\cos^2\theta}}.
\]
方向由 $E_r,E_\theta$ 共同确定；轴线上电场沿 $\vect p$，中垂面上电场与 $\vect p$ 反向。

电场线满足
\[
  \frac{\dd r}{r\dd\theta}=\frac{E_r}{E_\theta}=2\cot\theta.
\]
积分后得到
\[
  \boxed{r=C\sin^2\theta},
\]
其中 $C>0$ 标记不同的电场线。
\end{xesolution}

\Needspace{0.35\textheight}
\begin{origexamplebox}{例题 29.5【直线电四极子】}
三个共线点电荷依次为 $-q,+2q,-q$，相邻电荷间距均为 $l$。求 $r\gg l$ 处电场强度的
大小和方向，并求在该近似下的电场线方程。
\quadrupoleDiagram
\end{origexamplebox}

\begin{xesolution}
取 $+2q$ 所在点为原点，以三电荷连线为极轴。电荷系的总电荷与偶极矩都为零，因此远场的
最低阶非零项是四极项。利用
\[
  \frac{1}{|\vect r-\vect r'|}
  =\frac1r\sum_{n=0}^{\infty}\left(\frac{r'}r\right)^nP_n(\cos\theta)
\]
并保留 $n=2$ 项，得到
\[
  \varphi(r,\theta)\simeq
  \frac{kql^2}{r^3}(1-3\cos^2\theta).
\]
于是
\[
  \boxed{E_r=\frac{3kql^2}{r^4}(1-3\cos^2\theta)},
\]
\[
  \boxed{E_\theta=-\frac{6kql^2}{r^4}\sin\theta\cos\theta}.
\]
电场大小为
\[
  \boxed{E=\frac{3kql^2}{r^4}
  \sqrt{1-2\cos^2\theta+5\cos^4\theta}}.
\]
可见四极场按 $1/r^4$ 衰减，比偶极场的 $1/r^3$ 更快。

由电场线微分方程
\[
  \frac{\dd r}{r\dd\theta}=\frac{E_r}{E_\theta}
  =-\frac{1-3\cos^2\theta}{2\sin\theta\cos\theta}
\]
积分得
\[
  \boxed{r=C\sin\theta\sqrt{|\cos\theta|}}.
\]
上下半空间分别选取常量 $C$；$\cos\theta=0$ 的中垂面是场线结构的分界面。
\end{xesolution}

\Needspace{0.35\textheight}
\begin{origexamplebox}{例题 29.6【无穷大厚板】}
在 $y,z$ 两个方向无限延伸的厚板，仅在 $-d<x<d$ 内有均匀体电荷密度 $\rho>0$。
初始时，一个质量为 $m$、带电量 $+q$ 的点电荷在原点以速度 $v_0$ 沿 $+x$ 方向运动，
求该运动的周期。
\slabDiagram
\end{origexamplebox}

\begin{xesolution}
先由平面对称性求电场。取横跨中平面的柱形高斯面，可得
\[
  \vect E(x)=
  \begin{cases}
    \dfrac{\rho x}{\varepsilon_0}\uvec x, & |x|<d,\\[6pt]
    \dfrac{\rho d}{\varepsilon_0}\operatorname{sgn}(x)\uvec x, & |x|\ge d.
  \end{cases}
\]
对题目给定的 $q>0$，电场力始终把电荷推离中平面。在板内有
\[
  \ddot x=\frac{q\rho}{m\varepsilon_0}x,
\]
其解为双曲函数而非简谐函数；离开厚板后电荷仍受到向外的恒力。因此电荷不会返回，
\textbf{不存在运动周期}。可见题设电荷符号与“求周期”的要求相互矛盾。

若题意原本是带电量 $-q$ 的粒子，其中 $q>0$ 表示电荷量大小，则板内运动才是简谐运动：
\[
  \ddot x+\omega^2x=0,
  \qquad
  \omega=\sqrt{\frac{q\rho}{m\varepsilon_0}}.
\]
当 $v_0\le\omega d$ 时粒子始终留在厚板内，周期为
\[
  \boxed{T=\frac{2\pi}{\omega}
  =2\pi\sqrt{\frac{m\varepsilon_0}{q\rho}}}.
\]
当 $v_0>\omega d$ 时粒子会越出厚板，但外侧恒力仍指向厚板。令
$u=\sqrt{v_0^2-\omega^2d^2}$，完整周期为
\[
  \boxed{T=\frac{4}{\omega}\arcsin\frac{\omega d}{v_0}
  +\frac{4u}{\omega^2d}}.
\]
第一项是粒子四次穿过半块厚板所用的时间，第二项是两侧板外往返时间之和。
\end{xesolution}

\Needspace{0.35\textheight}
\begin{origexamplebox}{例题 29.7【电子透镜】}
电子显微镜中的一束电子先经电势差 $U_0$ 加速，再掠过一根与电子束原方向垂直的绷紧带电
细金属线。金属线带正电，线电荷密度为 $\lambda$；电场很弱，可把电子轨迹近似为直线。
若入射电子的瞄准参数为 $b$，求电子通过金属线后的偏转角 $\alpha$。
\electronLensDiagram
\end{origexamplebox}

\begin{xesolution}
电子经电势差 $U_0$ 加速后满足
\[
  \frac12m_ev^2=eU_0.
\]
取金属线为 $z$ 轴，未偏转轨迹为 $y=b$、$x=vt$。正线电荷在电子位置产生的电场为
\[
  \vect E=\frac{\lambda}{2\pi\varepsilon_0(x^2+b^2)}
  (x\uvec x+b\uvec y).
\]
电子带电 $-e$。由于关于 $x=0$ 对称，纵向冲量相消；横向冲量为
\begin{align*}
  \Delta p_y
  &=-\int_{-\infty}^{\infty}
  \frac{e\lambda b}{2\pi\varepsilon_0(x^2+b^2)}\frac{\dd x}{v}\\
  &=-\frac{e\lambda}{2\varepsilon_0v}.
\end{align*}
弱偏转条件下 $\alpha\simeq|\Delta p_y|/(m_ev)$，所以
\[
  \alpha=\frac{e\lambda}{2\varepsilon_0m_ev^2}
  =\boxed{\frac{\lambda}{4\varepsilon_0U_0}}.
\]
在无限长直线、弱偏转近似下，结果与瞄准参数 $b$ 无关；$b$ 只决定电子从金属线哪一侧通过。
正线电荷吸引电子，因此偏转方向指向金属线。若 $\alpha$ 不够小或金属线长度有限，则必须
使用真实弯曲轨迹和边缘场重新计算。
\end{xesolution}
